%%=============================================================================
%% ANEXO A
%%=============================================================================
\chapter{Anexo: Reproducibilidad de los experimentos}
\label{anexo:reproducibilidad}

Este anexo reúne la información necesaria para reproducir los experimentos del Capítulo~\ref{cap:resultados}.

\section{Repositorios}

\begin{table}[H]
\centering
\caption{Repositorios del proyecto}
\label{tab:repositorios}
\begin{tabularx}{\textwidth}{@{}L{0.20\textwidth}L{0.42\textwidth}Y@{}}
\toprule
\textbf{Repositorio} & \textbf{URL} & \textbf{Contenido} \\
\midrule
QPP-RNG-IoT & \url{https://github.com/QPP-IoT-Lib/QPP-RNG-IoT} & Generador, arnés, \textit{firmware} de pruebas, resultados \\
ASCON-IoT & \url{https://github.com/QPP-IoT-Lib/ASCON-IoT} & Ascon-AEAD128, vectores, validación y desempeño \\
QPP-IoT & \url{https://github.com/QPP-IoT-Lib/qpp-iot} & Fachada \texttt{QppEntropy} y fuentes por placa \\
\bottomrule
\end{tabularx}
\end{table}

Los resultados detallados, con las salidas completas de las herramientas del NIST, se encuentran en el directorio \texttt{hil-results} del repositorio QPP-RNG-IoT. Los archivos hexadecimales del \textit{firmware} usado en las capturas largas de AVR se distribuyen precompilados, con sus sumas SHA-256, en \texttt{hil-firmware/portable}.

\section{Comandos principales}

\paragraph{Pruebas y comparación en el anfitrión.}
\begin{lstlisting}
cargo test --workspace --release
cargo xtask compare
cargo run --release -p stats --bin stats-cli -- full \
    --dir <directorio-con-muestras> --out stats.json
\end{lstlisting}

\paragraph{Programación y captura en las placas} (placa: \texttt{mega2560}, \texttt{nano}, \texttt{nrf52840} o \texttt{esp32}; modo: \texttt{ref}, \texttt{rawsrc} o \texttt{iot}, y en la ESP32 también \texttt{rawsrc-noadc}). En la ESP32, el programa de captura reinicia la placa con \texttt{--reset-rts}.
\begin{lstlisting}
hil-firmware/tools/run.sh nano iot
uv run hil-firmware/tools/capture.py --port <puerto> \
    --records 2500 --out hil-results/raw/nano-iot
python3 hil-firmware/tools/extract.py iot \
    hil-results/raw/nano-iot.bin hil-results/samples/nano-iot.bin
hil-firmware/tools/smoke.sh nano
\end{lstlisting}

\paragraph{Estimación de min-entropía.}
\begin{lstlisting}
ea_non_iid -v <archivo> 8
ea_iid -v <archivo> 8
ea_non_iid -v -l 0,1000000 <residuos> 8    # E4, por ciclo
\end{lstlisting}

\paragraph{Traza por ciclo de la referencia en un anfitrión (E4).}
\begin{lstlisting}
cargo run --release -p qpp-rng-firmware --example pi_trace -- \
    <ciclos> <salida.bin>
\end{lstlisting}

\paragraph{Ascon-AEAD128.}
\begin{lstlisting}
cargo test --workspace --all-targets      # repositorio ASCON-IoT
scripts/build-avr.sh atmega328p           # repositorio QPP-IoT
\end{lstlisting}

\paragraph{Pila completa de extremo a extremo en las placas (repositorio QPP-IoT).}
\begin{lstlisting}
scripts/build-e2e.sh atmega328p           # AVR; grabar el ELF con avrdude
(cd firmware/e2e && cargo build --release --target thumbv7em-none-eabihf)
scripts/build-e2e.sh esp32                # ESP32; grabar con espflash
uv run scripts/e2e-check.py --port <puerto> --runs 3 \
    --ascon-c <ascon_c_driver>   # nRF52840: --reset break; ESP32: --reset rts
\end{lstlisting}

\section{Formato de los registros}

Cada \textit{firmware} de pruebas transmite por UART, a 115\,200 baudios, una línea ASCII de encabezado que empieza con \texttt{QPPHIL} y describe el modo y los parámetros (por ejemplo, \texttt{QPPHIL avr iot-wdt N=6 credit=1 samples\_per\_byte=8}), seguida de registros de 9 bytes: \texttt{0xA5} y dos enteros de 32 bits en orden \textit{little-endian}. El programa de captura descarta el byte de sincronización y guarda 8 bytes por registro; solo acepta un registro si el byte siguiente es el de sincronización del próximo, y cuenta los rechazos en el archivo de metadatos. En el modo de traza, los dos enteros son $n_p$ y $\Delta t$; en el modo de fuente cruda, la muestra de la fuente (y, en el nRF52840 y la ESP32, una palabra del RNG); en el modo IoT, cuatro bytes de salida y los ciclos que tomaron.

\section{Configuración de las herramientas del NIST}

La batería SP 800-22 del arnés copia el directorio \texttt{templates} de STS desde el directorio del ejecutable \texttt{assess}. Por eso, el directorio de STS debe estar en la ruta de búsqueda del sistema; no basta con un enlace simbólico al ejecutable. Antes de interpretar un informe, conviene comprobar que el archivo \texttt{stats.txt} de cada prueba no contenga la palabra \texttt{ABORTED}, y que el resultado de SP 800-22 en el JSON del arnés esté marcado como interpretado (\texttt{parsed}): si \texttt{assess} no está en la ruta de búsqueda, la batería puede reportar un veredicto global favorable sin haber ejecutado SP 800-22. La instalación usada en WSL se validó reproduciendo, fila por fila, el informe de SP 800-22 de la salida de 1\,MB del nRF52840 obtenido en macOS (1\,440 de 1\,452 combinaciones).
